% GENERATED FILE. Edit canonical lessons, then run npm run generate:books.
% canonical-source-hash: fnv1a64:54b73f59213820d7
% canonical-lessons: PA-C204-child-return, PA-C204-bachchian, PA-C204-bachchian-review, PA-C204-with-children, PA-C204-payoff

\chapter{With Children}
\label{ch:pa-a2-bachchian}
\hlchaptermodality{204}

\begin{chapteropening}
\textbf{By the end of this chapter:} \emph{I can recognise and use one children form before a familiar postposition.}

Recognise and produce one children form in a bounded familiar-postposition frame through listening, speaking, reading, and writing.
\end{chapteropening}

\section[\texorpdfstring{\pa{ਬੱਚਾ} (bachchaa)}{ਬੱਚਾ (bachchaa)}]{Return to the familiar child noun}

\label{lesson:PA-C204-child-return}



\begin{quote}
\emph{Hear:} \emph{bachchā}; YOU SAY: \textbf{\pa{ਬੱਚਾ}}
\end{quote}

\begin{grammarlens}[title={Grammar Lens}]
Keep this familiar singular noun ready. The next lesson will own one new form; this opener does not reveal it.
\end{grammarlens}

\subsection*{Your turn}
\begin{itemize}
\raggedright
  \item \emph{Hear:} \emph{a child}; YOU SAY: \textbf{\pa{ਬੱਚਾ}}
  \item Say the familiar noun once, then stop.
\end{itemize}

\begin{tcolorbox}[breakable,colback=teal!4,colframe=teal!35!black,title={Before you move on}]
Produce the familiar noun from meaning. This opener owns no new form.
\end{tcolorbox}
\section[\texorpdfstring{\pa{ਬੱਚਿਆਂ} (bachchiaan)}{ਬੱਚਿਆਂ (bachchiaan)}]{One whole form for children before a postposition}

\label{lesson:PA-C204-bachchian}



\begin{quote}
Retrieve the familiar singular noun for a child, then stop.
\end{quote}

\begin{grammarlens}[title={Grammar Lens: direct owner}]
\begin{quote}
\textbf{\pa{ਬੱਚਿਆਂ}} — \emph{bachchiā̃} — children, before a postposition
\end{quote}

Hear and say the whole form before reading it. Keep \textbf{\pa{ਬੱਚਿਆਂ}} whole. Use this form for “children” before the familiar \textbf{\pa{ਨਾਲ}}. The next lesson retrieves the new form alone before that old postposition returns.
\end{grammarlens}

\begin{cousinweb}
Punjabi University Patiala's noun grammar explains the form pattern used here. Its lexical statistics identifies this exact word as the plural-oblique form of the familiar child noun, and University Lesson 12 independently presents the same pattern. Sources: \href{https://www.learnpunjabi.org/noun.html}{Punjabi University noun grammar}, \href{https://www.learnpunjabi.org/statistics.html}{Punjabi University lexical statistics}, and \href{https://pt.learnpunjabi.org/lesson12.aspx?l=12}{Punjabi University Lesson 12}.

Chapter 55 records that the origin of \textbf{\pa{ਬੱਚਾ}} remains unsettled. This lesson makes no new historical claim and keeps the new form whole.
\end{cousinweb}

\subsection*{Your turn}
\begin{itemize}
\raggedright
  \item \emph{Hear:} \emph{bachchiā̃}; YOU SAY: \textbf{\pa{ਬੱਚਿਆਂ}}
  \item \emph{Read it:} \textbf{\pa{ਬੱਚਿਆਂ}}
  \item \emph{Write it:} copy \textbf{\pa{ਬੱਚਿਆਂ}} once, checking the whole form
\end{itemize}

\begin{tcolorbox}[breakable,colback=teal!4,colframe=teal!35!black,title={Before you move on}]
Produce the form for “children” used before a postposition.
\end{tcolorbox}
\section[\texorpdfstring{\pa{ਬੱਚਿਆਂ} (bachchiaan)}{ਬੱਚਿਆਂ (bachchiaan)}]{Retrieve the new children form on its own}

\label{lesson:PA-C204-bachchian-review}



\begin{quote}
\emph{Hear:} \emph{bachchiā̃}; YOU SAY: \textbf{\pa{ਬੱਚਿਆਂ}}
\end{quote}

\subsection*{Your turn}
\begin{itemize}
\raggedright
  \item \emph{Hear:} \emph{children, before a postposition}; YOU SAY: \textbf{\pa{ਬੱਚਿਆਂ}}
  \item \emph{Read it:} \textbf{\pa{ਬੱਚਿਆਂ}}
  \item \emph{Write it:} cover and write \textbf{\pa{ਬੱਚਿਆਂ}}
\end{itemize}

\begin{grammarlens}[title={Grammar Lens}]
Keep one whole form and one job here. The familiar postposition returns only after this separate retrieval.
\end{grammarlens}

\begin{tcolorbox}[breakable,colback=teal!4,colframe=teal!35!black,title={Before you move on}]
Say and then write the form once from memory.
\end{tcolorbox}
\section[\texorpdfstring{\pa{ਬੱਚਿਆਂ} \pa{ਨਾਲ} (bachchiaan naal)}{ਬੱਚਿਆਂ ਨਾਲ (bachchiaan naal)}]{Put the new form before one familiar postposition}

\label{lesson:PA-C204-with-children}



\begin{quote}
Retrieve the new children form, then the familiar word for “with.”
\end{quote}

\begin{grammarlens}[title={Grammar Lens}]
\begin{quote}
\textbf{\pa{ਬੱਚਿਆਂ} \pa{ਨਾਲ}} — \emph{bachchiā̃ nāl} — with children
\end{quote}

For this one frame, place the whole children form before \textbf{\pa{ਨਾਲ}}. The noun form comes first; the familiar postposition follows it. Nothing else is added.
\end{grammarlens}

\subsection*{Your turn}
\begin{itemize}
\raggedright
  \item \emph{Hear:} \emph{bachchiā̃ nāl}; YOU SAY: \textbf{\pa{ਬੱਚਿਆਂ} \pa{ਨਾਲ}}
  \item \emph{Read it:} \textbf{\pa{ਬੱਚਿਆਂ} \pa{ਨਾਲ}}
  \item \emph{Write it:} cover and write \textbf{\pa{ਬੱਚਿਆਂ} \pa{ਨਾਲ}}
\end{itemize}

\begin{tcolorbox}[breakable,colback=teal!4,colframe=teal!35!black,title={Before you move on}]
Produce the complete “with children” frame.
\end{tcolorbox}
\section[\texorpdfstring{\pa{ਬੱਚਿਆਂ} \pa{ਨਾਲ} (bachchiaan naal)}{ਬੱਚਿਆਂ ਨਾਲ (bachchiaan naal)}]{A bounded with-children payoff}

\label{lesson:PA-C204-payoff}



\begin{quote}
Retrieve the new children form alone before rebuilding its familiar postposition frame.
\end{quote}

\subsection*{Your turn: four skills}
\begin{itemize}
\raggedright
  \item \emph{Hear:} \emph{bachchiā̃ nāl}; YOU SAY: \textbf{\pa{ਬੱਚਿਆਂ} \pa{ਨਾਲ}}
  \item \emph{Read it:} \textbf{\pa{ਬੱਚਿਆਂ} \pa{ਨਾਲ}}
  \item \emph{Write it:} cover the line and write \textbf{\pa{ਬੱਚਿਆਂ} \pa{ਨਾਲ}}
\end{itemize}

\begin{tcolorbox}[breakable,colback=teal!4,colframe=teal!35!black,title={Before you move on}]
Say the new form, then produce the complete familiar-postposition frame.

This is the twenty-seventh prerequisite tranche exposed by Punjabi A2 mock drafting. It does not unblock Punjabi A2 mock 1 by itself. It is not an A2 mock, not A2 readiness or attainment, and not external or human validation.
\end{tcolorbox}
